% =================================================
% Часть 2. Решение логарифмического неравенства
% =================================================
\worksheettitle
  {Часть 2. Решаем логарифмическое неравенство}
  {От двух замен к точному ответу}
\studentnameblock

\begin{problembox}
  Решите неравенство
  \[
    \lg^4((x^2-4)^2)-\lg^2((x^2-4)^4)\ge192.
  \]
\end{problembox}

\begin{checkpointbox}[Вспомните ключевые переходы]
  Заполните по памяти. Если один из переходов вызывает затруднение,
  вернитесь к соответствующему разделу первой части.
  \[
    \lg^2 A=\answerblank[34mm],
    \qquad
    \lg(z^2)=\answerblank[42mm]\quad(z\ne0),
  \]
  \[
    0<|u|\le a
    \Longleftrightarrow
    \answerblank[108mm].
  \]
\end{checkpointbox}

\section{ОДЗ и преобразование логарифмов}

Аргументы логарифмов должны быть положительными:
\[
  (x^2-4)^2>0
  \quad\Longleftrightarrow\quad
  x\ne\answerblank[14mm],\quad x\ne\answerblank[14mm].
\]

На ОДЗ:
\[
  \lg((x^2-4)^2)=2\lg\answerblank[30mm],
\]
\[
  \lg((x^2-4)^4)=4\lg\answerblank[30mm].
\]
\newpage
\begin{reflectionbox}[Контроль равносильности]
  Почему запись \(4\lg(x^2-4)\) потеряла бы часть допустимых значений?
  \writinglines{2}
\end{reflectionbox}

\begin{fillbox}[Переписываем исходное неравенство целиком]
  Важно сохранить внешние степени: в четвёртую степень возводится значение
  первого логарифма, а в квадрат --- значение второго.
  \[
    \begin{aligned}
      &\bigl[\lg((x^2-4)^2)\bigr]^4
      -\bigl[\lg((x^2-4)^4)\bigr]^2\ge192\\
      &\quad\Longleftrightarrow
      \bigl[2\lg|x^2-4|\bigr]^4
      -\bigl[4\lg|x^2-4|\bigr]^2\ge192.
    \end{aligned}
  \]
  Какое выражение повторяется в обеих частях полученного неравенства?
  \[
    \answerblank[65mm].
  \]
\end{fillbox}

\section{Сводим задачу к квадратному неравенству}

Положим
\[
  y=\lg|x^2-4|.
\]
Тогда
\[
  \begin{aligned}
    (2y)^4-(4y)^2&\ge192,\\
    16y^4-16y^2&\ge192,\\
    y^4-y^2-12&\ge0.
  \end{aligned}
\]

Введите замену и ограничение:
\[
  t=\answerblank[20mm],
  \qquad t\answerblank[10mm]0.
\]
Решите получившееся квадратное неравенство:
\[
  t^2-t-12=(t-\answerblank[12mm])(t+\answerblank[12mm])\ge0,
\]
\[
  t\in\answerblank[75mm].
\]
С учётом ограничения на \(t\):
\[
  t\ge\answerblank[12mm].
\]

\begin{checkpointbox}[Не потеряйте отрицательную ветвь]
  \[
    \begin{aligned}
      t\ge4
      &\Longrightarrow y^2\ge4
      \Longleftrightarrow |y|\ge\answerblank[10mm],\\
      &\Longleftrightarrow \answerblank[75mm].
    \end{aligned}
  \]
\end{checkpointbox}

\section{Выполняем обратный переход}

Так как \(y=\lg|x^2-4|\), получаем:
\[
  \lg|x^2-4|\le-2
  \qquad\text{или}\qquad
  \lg|x^2-4|\ge2.
\]
Десятичный логарифм возрастает, поэтому:
\[
  0<|x^2-4|\le\answerblank[22mm]
  \qquad\text{или}\qquad
  |x^2-4|\ge\answerblank[22mm].
\]

\argregimesstudentfigure

\newpage
\section{Разбираем два случая}

\begin{fillbox}[Ветвь малого аргумента]
  Используйте правило из первой части:
  \[
    0<|x^2-4|\le\frac1{100}.
  \]
  После раскрытия условия:
  \[
    -\frac1{100}\le x^2-4<0
    \quad\text{или}\quad
    0<x^2-4\le\frac1{100}.
  \]
  После прибавления \(4\):
  \[
    \frac{\phantom{000}}{100}\le x^2<4
    \quad\text{или}\quad
    4<x^2\le\frac{\phantom{000}}{100}.
  \]
  Запишите четыре промежутка для \(x\):
  \workarea{42mm}
\end{fillbox}

\begin{fillbox}[Ветвь большого аргумента]
  \[
    |x^2-4|\ge100
    \Longleftrightarrow
    x^2-4\le\answerblank[18mm]
    \quad\text{или}\quad
    x^2-4\ge\answerblank[18mm].
  \]
  Какая ветвь невозможна и почему?
  \writinglines{2}
  Оставшаяся ветвь даёт:
  \[
    x^2\ge\answerblank[16mm]
    \Longleftrightarrow
    |x|\ge\answerblank[22mm].
  \]
  \[
    x\in\answerblank[100mm].
  \]
\end{fillbox}

\section{Собираем ответ}

\finalnumberlinestudentfigure

\begin{resultbox}[Точный ответ]
  \writinglines{4}
\end{resultbox}

\begin{checkpointbox}[Проверка характерных точек]
  Проверьте подстановкой или по полученным условиям:
  \[
    x=0,\qquad x=2,\qquad
    x=\frac{\sqrt{399}}{10},\qquad x=2\sqrt{26}.
  \]
  Какая точка не принадлежит ОДЗ? Какие две являются граничными решениями?
  \writinglines{2}
\end{checkpointbox}

\newpage
\section{Самостоятельно записываем решение}

\begin{practicebox}[Самостоятельная запись]
  Закройте предыдущие страницы и запишите связное решение. Проверьте полноту
  только после завершения:
  \begin{enumerate}
    \item ОДЗ и формулы с модулем;
    \item замены \(y=\lg|x^2-4|\) и \(t=y^2\ge0\);
    \item обе ветви обратного перехода;
    \item решение двух условий относительно \(x\);
    \item точный ответ.
  \end{enumerate}
  \workarea{150mm}
\end{practicebox}
